\documentclass[10.5pt,a4paper]{article}
\usepackage{fontspec}
\usepackage{xeCJK}
\usepackage[margin=18mm]{geometry}
\usepackage{xcolor}
\usepackage{array}
\usepackage{longtable}
\usepackage{booktabs}
\usepackage{fancyhdr}
\usepackage{hyperref}
\setmainfont{Arial}
\setCJKmainfont{SimHei}
\hypersetup{colorlinks=true,urlcolor=blue}
\pagestyle{fancy}\fancyhf{}\lhead{Evidence-aware Retrieval Study}\rhead{TR-2026-08}\cfoot{\thepage}
\setlength{\parindent}{0pt}\setlength{\parskip}{4pt}
\begin{document}
{\fontsize{21}{25}\selectfont\bfseries 面向企业知识库的证据保留式文档解析研究}\par
{\color{gray}\fontsize{12}{15}\selectfont Evidence-preserving Document Parsing for Enterprise Knowledge Bases}\par\vspace{3mm}{\color{blue}\rule{\textwidth}{1.8pt}}\par\vspace{5mm}
\begin{tabular}{p{31mm}p{55mm}p{31mm}p{55mm}}
作者 & 陈琳，周海，王倩 & 文档类型 & 技术研究报告\\
版本 & 1.2 & 发布日期 & 2026-08-17\\
数据集 & Shared Dev v1 & 状态 & Reviewed\\
\end{tabular}
\section*{摘要}
企业知识库中的 PDF、扫描件和 Office 文档同时包含正文、表格、图注和引用。只抽取纯文本会损失阅读顺序和来源证据。本文提出一套保留页面、边界框、表格网格和引用关系的解析流程，并在合成但业务真实的文档集上进行评估。结果表明，结构证据能够减少表格字段错绑和引用目标误判。

\textbf{关键词：} 文档解析；证据保留；表格绑定；标题树；引用关系；质量门
\newpage
\section{引言}
\subsection{背景}
企业文档通常由多个团队维护，标题、版本和附件并不统一。检索系统如果只保存 Markdown，会难以回答“这个数字来自哪一页”或“该结论引用了哪篇文献”。已有研究指出，布局与语义结构对文档问答同样重要 [1-2]。
\subsection{问题定义}
给定一个包含文本块、表格、图片和引用的文档，系统需要输出统一的 ParsedDocument，并保留可回溯的 source locator。质量层还要区分 verified、inferred 和 manual review required，而不是把所有推断都当成事实 [3]。
\subsection{贡献}
\begin{enumerate}
\item 定义包含页码、bbox、表格 cell/span 和原生产物引用的证据模型；
\item 提出针对多级表头和跨页表格的安全绑定规则；
\item 建立标题树和正文引用到参考文献的可验证关系；
\item 使用固定样例和质量门报告评估安全降级行为。
\end{enumerate}
\section{方法}
\subsection{输入预检}
预检器记录文件类型、页数、文本层、扫描页比例和语言提示。路由器根据能力声明选择解析器，解析器实际执行来源写入 provenance [4]。
\subsection{统一结构}
每个文本块有稳定 block id、阅读顺序和 source anchor。表格 cell 保留 start row、start col、row span、col span；合并表头生成完整 column path，例如“训练集 / 准确率”。
\subsection{质量门}
质量门按照以下规则决定状态：证据充分且关系唯一时为 verified；只有顺序和语义支持时为 inferred；目标不唯一时进入 manual review；关键事实缺失且其他模型可补足时进入 reparse required [5]。
\newpage
\section{实验设置}
\subsection{样例与指标}
\begin{tabular}{p{45mm}c c p{62mm}}
\toprule
能力 & 正例数 & 反例数 & 主要指标\\\midrule
普通正文与阅读顺序 & 2 & 1 & 文本完整度、block 顺序\\
多级表头与字段绑定 & 2 & 2 & column path、row key、来源\\
标题层级恢复 & 2 & 1 & parent\_child 唯一性\\
编号引用绑定 & 2 & 2 & reference\_of 目标唯一性\\
OCR 图片质量 & 2 & 2 & token 可读性、低质量拦截\\
\bottomrule
\end{tabular}
\subsection{评估流程}
每个样例都使用相同的输入摘要和 manifest commit。运行记录包含路由决定、实际解析器、耗时、ParsedDocument schema、质量状态、问题数量和产物位置。失败和不可用结果不会从批量报告中删除。
\subsection{表格结果}
\begin{tabular}{p{40mm}c c c c}
\toprule
样例 & 网格恢复 & 表头路径 & 字段绑定 & 最终状态\\\midrule
采购比价正例 & pass & 完整 & verified & pass\\
跨页续表正例 & pass & 完整 & verified & pass with warnings\\
列漂移反例 & partial & 不确定 & unavailable & manual review\\
无稳定行键 & partial & 表级 & unavailable & manual review\\
\bottomrule
\end{tabular}
\section{讨论}
\subsection{安全降级}
在表格结构不唯一时，减少绑定数量比输出错误关系更安全。质量报告必须同时记录缺失证据、受影响 block 和重新解析建议 [6]。
\subsection{局限性}
合成样例不能覆盖所有扫描噪声和供应商格式。后续 Acceptance 子集应加入未参与规则调试的公开文档，并由人工标注确认。
\newpage
\section{结论}
本文展示了从预检、解析、统一输出到质量门的完整证据链。实验说明，多级标题、表格网格和引用关系需要共同的来源模型，不能由质量层凭空补造。对于无法唯一判断的结构，系统应明确交给人工复核或重新解析。

\section*{参考文献}
\begin{enumerate}
\item Li, M.; Zhang, Y. Document layout analysis for retrieval-augmented generation[J]. Journal of Information Processing, 2024, 32(4): 201-218. DOI:10.1000/jip.2024.03204.
\item 王伟, 陈晨. 面向复杂版面的文档结构恢复[J]. 中文信息学报, 2023, 37(6): 45-58. DOI:10.1000/cip.2023.3706.
\item Pujara, J.; Singh, H. Evidence graphs for trustworthy document QA[C]. Proceedings of DocEng, 2024: 55-63. DOI:10.1000/doceng.2024.006.
\item Liu, Q.; Chen, R. OCR confidence calibration in enterprise archives[J]. Pattern Recognition Letters, 2022, 158: 12-20. DOI:10.1000/prl.2022.15802.
\item 赵敏, 周海. 表格字段绑定的来源约束与安全降级[J]. 软件学报, 2025, 36(8): 1771-1787. DOI:10.1000/js.2025.3608.
\item Smith, A.; Patel, R. Human review gates for semantic extraction[C]. Proceedings of ICDAR, 2023: 410-419. DOI:10.1000/icdar.2023.041.
\end{enumerate}
\vspace{6mm}
\textbf{引用规则：} 正文中的 [1-2] 指向参考文献 1 和 2；[3]、[4]、[5]、[6] 分别指向对应编号。\par
\vfill\color{gray}\small Synthetic technical report; authors, measurements and DOI values are fictional but internally consistent.
\end{document}